% Part: computability
% Chapter: computability-theory
% Section: halting-problem

\documentclass[../../../include/open-logic-section]{subfiles}

\begin{document}

\olfileid{cmp}{thy}{hlt}
\olsection{مشكلة التوقف}

بحكم البناء، تكون الدالة الجزئية الشاملة القابلة للحساب
~$\fn{Un}(e,x)$ معرفة إذا وفقط إذا أنتج حساب الدالة التي يرمز إليها
~$e$ قيمة عند الدخل~$x$. ومن الطبيعي أن نسأل هل يمكننا تقرير ما إذا
كان الأمر كذلك. والواقع أنه لا يمكن. وهذا يعني، في نموذج الحساب بآلة
تورنغ، أن مسألة توقف آلة تورنغ معطاة عند دخل معطى غير قابلة للقرار
حسابيًا. ولذلك تعرف المبرهنة الآتية باسم «عدم قابلية مشكلة التوقف
للقرار». وسنقدم أدناه برهانين؛ يواصل الأول مسار مناقشتنا السابقة، أما
الثاني فأكثر مباشرة.

\begin{thm}
\ollabel{thm:halting-problem}
لتكن
\[
h(e, x)  =
\begin{cases}
1 & \text{إذا كانت\/ $\fn{Un}(e, x)$ معرفة} \\
0 & \text{في غير ذلك.}
\end{cases}
\]
عندئذ لا تكون $h$ قابلة للحساب.
\end{thm}

\begin{proof}
افترض أن $h$ قابلة للحساب. سنبين أن هذا يتيح لنا تعريف دالة كلية
شاملة قابلة للحساب. عرف
\[
\fn{Un'}(e,x) =
\begin{cases}
\fn{Un}(e,x) & \text{إذا كان $h(e,x) = 1$} \\
0 & \text{في غير ذلك.}
\end{cases}
\]
لكن $\fn{Un'}(e,x)$ الآن دالة كلية، وتكون قابلة للحساب إذا كانت~$h$
كذلك. فمثلًا يمكننا تعريف $g$ بالعودية البدائية هكذا:
\begin{align*}
g(0, e, x) & \simeq 0 \\
g(y+1, e, x) & \simeq \fn{Un}(e,x);
\end{align*}
وعندئذ
\[
\fn{Un'}(e,x) \simeq g(h(e,x),e,x).
\]
ولما كانت $\fn{Un'}(e,x)$ تطابق $\fn{Un}(e,x)$ حيثما كانت الأخيرة
معرفة، فإن $\fn{Un'}$ شاملة للدوال الجزئية القابلة للحساب التي تصادف
أن تكون كلية. لكن هذا يناقض \olref[nou]{thm:no-univ}.
\end{proof}

\begin{proof}
افترض أن $h(e,x)$ قابلة للحساب. عرف الدالة $g$ بـ
\[
g(x) =
\begin{cases}
  0                & \text{إذا كان $h(x,x) = 0$} \\
  \fundefined & \text{في غير ذلك.}
\end{cases}
\]
الدالة $g$ جزئية قابلة للحساب. فمثلًا يمكن تعريفها بأنها
$\umin{y}{h(x,x)=0}$. ولذلك يوجد~$e$ بحيث
$g(x)\simeq\fn{Un}(e,x)$ لكل~$x$. فهل $g$ معرفة عند $e$؟ إذا كانت
معرفة، لكان $h(e,e)=0$ بحسب تعريف~$g$ (إذ لا يمكن أن تكون~$g$ معرفة
إلا إذا كان $h(e,e)=0$). ويعني هذا، بحسب تعريف~$h$، أن $\fn{Un}(e,e)$ غير
معرفة. ومن افتراضنا أن $g(x)\simeq\fn{Un}(e,x)$ لكل~$x$، نحصل على أن
$g(e)$ غير معرفة، وهذا تناقض. ومن جهة أخرى، إذا كانت $g(e)$ غير
معرفة، كان $h(e,e)\neq0$، ومن ثم $h(e,e)=1$. ويترتب على ذلك أن
$\fn{Un}(e,e)$ معرفة. لكن لما كان $g(x)\simeq\fn{Un}(e,x)$، لكانت
$g(e)$ معرفة أيضًا. وهذا تناقض مرة أخرى.
\end{proof}

\begin{tagblock}{TMs}
\begin{explain}
يمكننا وصف هذه الحجة بدلالة آلات تورنغ. افترض وجود آلة تورنغ~$H$ تأخذ
دخلًا هو وصف آلة تورنغ~$E$ ودخلًا~$x$، وتقرر ما إذا كانت $E$ تتوقف
عند الدخل~$x$. عندئذ يمكننا بناء آلة تورنغ أخرى~$G$ تأخذ دخلًا
واحدًا~$x$، وتشغل $H$ لتقرر ما إذا كانت الآلة $M_x$ ذات الفهرس~$x$
تتوقف عند الدخل~$x$، ثم تفعل العكس. وبعبارة أخرى، إذا أفادت $H$ بأن
$M_x$ تتوقف عند الدخل~$x$، دخلت $G$ في حلقة لا نهائية؛ وإذا أفادت
$H$ بأن $M_x$ لا تتوقف عند الدخل~$x$، توقفت $G$ فحسب. فهل تتوقف $G$
حين يكون فهرسها نفسه دخلها؟ تبين الحجة أعلاه أنها تتوقف إذا وفقط إذا
لم تتوقف، وهذا تناقض. فلا بد أن يكون افتراض وجود آلة تورنغ~$H$ كهذه
باطلًا.
\end{explain}
\end{tagblock}

\end{document}
